% ============================================================
% PROPOSAL PROYEK AKHIR — FARCHAN DEANO MUHAMMAD (CAN)
% PENGEMBANGAN SIMULASI INTERAKTIF BERLEVEL DENGAN MEKANISME
% HUMAN-IN-THE-LOOP PADA SISTEM LITERASI KECERDASAN BUATAN
% UNTUK SISWA SMP
%
% Format: PENS Standard (A4, 12pt TNR, 1.5 spacing, 4cm/3cm margin)
% Mengikuti pola format Izzah (PilAItes 2025)
% Skeleton file — empty sections siap diisi Phase 2
% ============================================================

\documentclass[12pt,a4paper,oneside,openany]{report}

% ============================================================
% PACKAGES
% ============================================================
\usepackage[top=4cm,left=4cm,right=3cm,bottom=3cm]{geometry}
\usepackage{graphicx}
\usepackage{booktabs}
\usepackage{array}
\usepackage{tabularx}
\usepackage{longtable}
\usepackage{float}
\usepackage{caption}
\usepackage{titlesec}
\usepackage{indentfirst}
\usepackage{setspace}
\usepackage{etoolbox}
\usepackage{fancyhdr}
\usepackage{afterpage}
\usepackage{amsmath}
\usepackage{amssymb}
\usepackage{url}
\usepackage{hyperref}
\hypersetup{colorlinks=false, pdfborder={0 0 0}}

% Times New Roman (pdflatex standard)
\usepackage{times}
\renewcommand{\rmdefault}{ptm}

% ============================================================
% LINE SPACING 1.5
% ============================================================
\setstretch{1.5}

% ============================================================
% PARAGRAPH INDENT (1.25 cm, PENS standard)
% ============================================================
\setlength{\parindent}{1.25cm}
\setlength{\parskip}{0pt}

% ============================================================
% CHAPTER / SECTION / SUBSECTION FORMATTING
% ============================================================
\titleformat{\chapter}[display]
  {\normalfont\fontsize{14pt}{16.8pt}\selectfont\bfseries\centering}
  {BAB \thechapter}
  {12pt}
  {\MakeUppercase}

\titleformat{\section}
  {\normalfont\fontsize{12pt}{14.4pt}\selectfont\bfseries}
  {\thesection}
  {1em}
  {\MakeUppercase}

\titleformat{\subsection}
  {\normalfont\fontsize{12pt}{14.4pt}\selectfont\bfseries}
  {\thesubsection}
  {1em}
  {}

\titleformat{\subsubsection}
  {\normalfont\fontsize{12pt}{14.4pt}\selectfont\bfseries\itshape}
  {\thesubsubsection}
  {1em}
  {}

\titlespacing*{\chapter}{0pt}{0pt}{24pt}
\titlespacing*{\section}{0pt}{18pt}{12pt}
\titlespacing*{\subsection}{0pt}{14pt}{8pt}

% ============================================================
% CAPTION FORMATTING (Pola Izzah)
% - Tabel: caption DI ATAS, font 11pt, label bold
% - Gambar: caption DI BAWAH, font 11pt, label bold
% ============================================================
\captionsetup[table]{font=small,labelfont=bf,name={Tabel},labelsep=period,justification=centering,aboveskip=6pt,belowskip=0pt,position=top}
\captionsetup[figure]{font=small,labelfont=bf,name={Gambar},labelsep=period,justification=centering,aboveskip=0pt,belowskip=6pt,position=bottom}

% ============================================================
% HEADER / FOOTER (page number top-right)
% ============================================================
\pagestyle{fancy}
\fancyhf{}
\fancyhead[R]{\thepage}
\renewcommand{\headrulewidth}{0pt}
\renewcommand{\footrulewidth}{0pt}

\fancypagestyle{plain}{
  \fancyhf{}
  \fancyhead[R]{\thepage}
  \renewcommand{\headrulewidth}{0pt}
}

% ============================================================
% BLANK PAGE COMMAND
% ============================================================
\newcommand{\blankpage}{%
  \afterpage{%
    \null
    \vfill
    \begin{center}
      {\small Halaman ini sengaja dikosongkan}
    \end{center}
    \vfill
    \thispagestyle{plain}
    \newpage
  }%
}

% ============================================================
% CUSTOM COMMANDS
% ============================================================
% \gbr{filename}{caption}{label}
\newcommand{\gbr}[3]{%
  \begin{figure}[H]
    \centering
    \includegraphics[width=\textwidth]{#1}
    \caption{#2}
    \label{#3}
  \end{figure}
}

% \gbrsrc{filename}{caption}{label}{source}
\newcommand{\gbrsrc}[4]{%
  \begin{figure}[H]
    \centering
    \includegraphics[width=\textwidth]{#1}
    \caption{#2}
    \label{#3}
    \begin{center}\small (Sumber: #4)\end{center}
  \end{figure}
}

% \tabel{caption}{content}{label}
\newcommand{\tabel}[3]{%
  \begin{table}[H]
    \centering
    \caption{#1}
    \label{#3}
    #2
  \end{table}
}

% ============================================================
% TOC FORMATTING
% ============================================================
\renewcommand{\contentsname}{DAFTAR ISI}
\renewcommand{\listfigurename}{DAFTAR GAMBAR}
\renewcommand{\listtablename}{DAFTAR TABEL}

% ============================================================
\begin{document}
% ============================================================

% ============================================================
% 1. COVER PAGE
% ============================================================
\begin{titlepage}
\centering
\vspace*{1cm}

{\fontsize{14pt}{16pt}\selectfont PROPOSAL PROYEK AKHIR}

\vspace{2cm}

{\fontsize{14pt}{18pt}\selectfont\bfseries\MakeUppercase{Pengembangan Simulasi Interaktif Berlevel dengan Mekanisme Human-in-the-Loop pada Sistem Literasi Kecerdasan Buatan untuk Siswa SMP}}

\vspace{3cm}

{\fontsize{12pt}{14pt}\selectfont Oleh:}\\[0.3cm]
{\fontsize{12pt}{14pt}\selectfont\bfseries Farchan Deano Muhammad}\\[0.2cm]
{\fontsize{12pt}{14pt}\selectfont NRP. 53226000xx}

\vspace{3cm}

{\fontsize{12pt}{14pt}\selectfont Dosen Pembimbing:}\\[0.3cm]
{\fontsize{12pt}{14pt}\selectfont [Nama Dosen Pembimbing]}\\[0.2cm]
{\fontsize{12pt}{14pt}\selectfont NIP. [NIP]}

\vspace{3cm}

{\fontsize{12pt}{14pt}\selectfont PROGRAM STUDI SARJANA TERAPAN TEKNOLOGI REKAYASA MULTIMEDIA}\\[0.2cm]
{\fontsize{12pt}{14pt}\selectfont JURUSAN TEKNOLOGI MULTIMEDIA DAN JARINGAN}\\[0.2cm]
{\fontsize{12pt}{14pt}\selectfont POLITEKNIK ELEKTRONIKA NEGERI SURABAYA}\\[0.2cm]
{\fontsize{12pt}{14pt}\selectfont 2025}

\end{titlepage}

% ============================================================
% 2. BLANK PAGE
% ============================================================
\blankpage

% ============================================================
% 3. HALAMAN PENGESAHAN
% ============================================================
\chapter*{HALAMAN PENGESAHAN}
\addcontentsline{toc}{chapter}{HALAMAN PENGESAHAN}

Proposal Proyek Akhir ini disusun dan diajukan oleh:

\vspace{0.4cm}

\begin{tabular}{ll}
Nama & : Farchan Deano Muhammad \\
NRP & : 53226000xx \\
Program Studi & : Sarjana Terapan Teknologi Rekayasa Multimedia \\
Jurusan & : Teknologi Multimedia dan Jaringan \\
\end{tabular}

\vspace{0.5cm}

Telah diperiksa dan disahkan oleh Dosen Pembimbing pada tanggal: \underline{\hspace{3cm}}

\vspace{1.5cm}

\begin{tabular}{ll}
Dosen Pembimbing, & \\
\\
\\
\\
\underline{\hspace{5cm}} & \\
[Nama Dosen Pembimbing] & \\
NIP. \underline{\hspace{3cm}} & \\
\end{tabular}

\blankpage

% ============================================================
% 4. ABSTRAK
% ============================================================
\chapter*{ABSTRAK}
\addcontentsline{toc}{chapter}{ABSTRAK}

% TODO Phase 2: Isi abstrak 200-250 kata + kata kunci
% Template: konteks → masalah → solusi → metode → hasil yang diharapkan
%
% [ABSTRAK — ISI DI SINI]

\vspace{0.5cm}

\noindent\textbf{Kata Kunci}: % [kata kunci 1, kata kunci 2, ...]

\blankpage

% ============================================================
% 5. ABSTRACT (English)
% ============================================================
\chapter*{ABSTRACT}
\addcontentsline{toc}{chapter}{ABSTRACT}

% TODO Phase 2: English version of abstrak
% [ABSTRACT — ISI DI SINI]

\vspace{0.5cm}

\noindent\textbf{Keywords}: % [keywords in English]

\blankpage

% ============================================================
% 6. DAFTAR ISI
% ============================================================
\chapter*{DAFTAR ISI}
\addcontentsline{toc}{chapter}{DAFTAR ISI}

\tableofcontents
\clearpage

% ============================================================
% 7. DAFTAR GAMBAR
% ============================================================
\chapter*{DAFTAR GAMBAR}
\addcontentsline{toc}{chapter}{DAFTAR GAMBAR}

\listoffigures
\clearpage

% ============================================================
% 8. DAFTAR TABEL
% ============================================================
\chapter*{DAFTAR TABEL}
\addcontentsline{toc}{chapter}{DAFTAR TABEL}

\listoftables
\clearpage

% ============================================================
% BAB 1 — PENDAHULUAN
% (UNCHANGED from existing proposal — JANGAN diubah, hanya format convert)
% Total target: ~5 halaman
% ============================================================
\chapter{PENDAHULUAN}
\label{chap:bab1}

% ------------------------------------------------------------
\section{LATAR BELAKANG}
\label{sec:latar-belakang}
% ------------------------------------------------------------
% TODO Phase 2: Copy Bab 1.1 dari proposal-can.tex existing
% Estimasi: ~2 halaman, 5-6 paragraf
% [ISI DI SINI]

% ------------------------------------------------------------
\section{PERMASALAHAN}
\label{sec:permasalahan}
% ------------------------------------------------------------
% TODO Phase 2: Copy Bab 1.2 dari proposal-can.tex existing
% Estimasi: ~0.5 halaman, 1 paragraf
% [ISI DI SINI]

% ------------------------------------------------------------
\section{BATASAN MASALAH}
\label{sec:batasan-masalah}
% ------------------------------------------------------------
% TODO Phase 2: Copy Bab 1.3 dari proposal-can.tex existing
% Estimasi: ~0.5 halaman, bullet list 4-5 poin
% [ISI DI SINI]

% ------------------------------------------------------------
\section{TUJUAN}
\label{sec:tujuan}
% ------------------------------------------------------------
% TODO Phase 2: Copy Bab 1.4 dari proposal-can.tex existing
% Estimasi: ~0.5 halaman
% [ISI DI SINI]

% ------------------------------------------------------------
\section{MANFAAT}
\label{sec:manfaat}
% ------------------------------------------------------------
% TODO Phase 2: Copy Bab 1.5 dari proposal-can.tex existing
% Estimasi: ~0.5 halaman, bullet 4 poin
% [ISI DI SINI]

% ------------------------------------------------------------
\section{SISTEMATIKA PENULISAN}
\label{sec:sistematika-penulisan}
% ------------------------------------------------------------
% TODO Phase 2: Copy Bab 1.6 dari proposal-can.tex existing
% Estimasi: ~0.5 halaman, paragraf per BAB
% [ISI DI SINI]

% ============================================================
% BAB 2 — KAJIAN PUSTAKA
% BEDA Can vs Dias — Can fokus Frontend/HITL/UI-UX
% Total target: ~8-10 halaman
% Pola Izzah: ringkas, detail di Bab 3
% ============================================================
\chapter{KAJIAN PUSTAKA}
\label{chap:bab2}

% Paragraf pembuka bab — TODO Phase 2
% [PARAGRAF PEMBUKA BAB 2]

% ------------------------------------------------------------
\section{DESKRIPSI PERMASALAHAN}
\label{sec:deskripsi-permasalahan}
% ------------------------------------------------------------
% TODO Phase 2: Copy dari proposal-can.tex existing + ringkas
% Estimasi: ~1 halaman, 3 paragraf
% [ISI DI SINI]

% ------------------------------------------------------------
\section{TEORI PENUNJANG}
\label{sec:teori-penunjang}
% ------------------------------------------------------------
% TODO Phase 2: Copy dari proposal-can.tex existing + ringkas (detail di Bab 3)
% Estimasi: ~5-6 halaman total, 0.5-1 hal per teori
% Format pola Izzah: definisi → penjelasan → posisi dalam penelitian

\subsection{Literasi Kecerdasan Buatan}
\label{subsec:literasi-ai}
% [ISI DI SINI]

\subsection{Human-in-the-Loop}
\label{subsec:hitl}
% [ISI DI SINI]

\subsection{Explainable Interface dan Confidence Score}
\label{subsec:xai-confidence}
% [ISI DI SINI]

\subsection{Game-Based Learning dan Scaffolding}
\label{subsec:gbl-scaffolding}
% [ISI DI SINI]

\subsection{Pedagogical Agent}
\label{subsec:pedagogical-agent}
% [ISI DI SINI]

\subsection{Child-Computer Interaction untuk Siswa SMP}
\label{subsec:cci}
% [ISI DI SINI]

\subsection{System Usability Scale}
\label{subsec:sus}
% [ISI DI SINI]

% ------------------------------------------------------------
\section{PENELITIAN TERKAIT}
\label{sec:penelitian-terkait}
% ------------------------------------------------------------
% TODO Phase 2: 4 penelitian terkait + Tabel State of the Art
% Estimasi: ~2-3 halaman

\subsection{AI Literacy in K-12}
\label{subsec:penelitian-1}
% [ISI DI SINI]

\subsection{Human-in-the-Loop Machine Learning}
\label{subsec:penelitian-2}
% [ISI DI SINI]

\subsection{Game-Based Learning in Computer Science}
\label{subsec:penelitian-3}
% [ISI DI SINI]

\subsection{Real-Time Hand Gesture Monitoring}
\label{subsec:penelitian-4}
% [ISI DI SINI — NOTE: revisi karena finger tracking login DEPRECATED]

\subsection{State of the Art}
\label{subsec:state-of-the-art}
% TODO Phase 2: Tabel 2.1 State of the Art (POLosan, format Izzah)
% Kolom: No | Penelitian | Metode/Fokus | Temuan | Celah yang Diisi
% [ISI DI SINI + Tabel 2.1]

% ============================================================
% BAB 3 — DESKRIPSI SISTEM
% BEDA Can vs Dias — DETAIL, scope Frontend/HITL/UI-UX
% Total target: ~18-22 halaman
% Urutan pola Izzah: Deskripsi Solusi → Metodologi → Desain Sistem → Jadwal
% ============================================================
\chapter{DESKRIPSI SISTEM}
\label{chap:bab3}

% Paragraf pembuka bab — TODO Phase 2
% [PARAGRAF PEMBUKA BAB 3]

% ------------------------------------------------------------
\section{DESKRIPSI SOLUSI}
\label{sec:deskripsi-solusi}
% ------------------------------------------------------------
% TODO Phase 2: Copy dari proposal-can.tex existing Bab 3.1
% Estimasi: ~1-2 halaman
% Konten: 4 prinsip akademis + 3 level progresi + Momo + tech stack ringkas
% [ISI DI SINI]

% ------------------------------------------------------------
\section{METODOLOGI PENELITIAN}
\label{sec:metodologi-penelitian}
% ------------------------------------------------------------
% TODO Phase 2: PINDAH dari posisi akhir (existing 3.4) ke sini (3.2)
% Estimasi: ~2-3 halaman
% Konten WAJIB: Waterfall (SDLC) 5 tahap + Fishbone 6M + SUS + Black-Box
% [ISI DI SINI]

\subsection{Pendekatan dan Jenis Penelitian}
\label{subsec:pendekatan-penelitian}
% [ISI DI SINI — Mixed Methods + R&D]

\subsection{Metode Pengembangan}
\label{subsec:metode-pengembangan}
% [ISI DI SINI — SDLC Waterfall 5 tahap]

\subsection{Analisis Masalah (Fishbone 6M)}
\label{subsec:analisis-masalah}
% [ISI DI SINI — Fishbone 6M: Man, Machine, Method, Material, Measurement, Environment]
% TODO: Tambahkan Gambar Fishbone Diagram

\subsection{Teknik Pengumpulan Data}
\label{subsec:teknik-pengumpulan-data}
% [ISI DI SINI — Data Log + Kuesioner SUS + Observasi]

\subsection{Teknik Analisis Data}
\label{subsec:teknik-analisis-data}
% [ISI DI SINI — Statistik Deskriptif + Fishbone + Thematic Analysis]

\subsection{Metode Pengujian}
\label{subsec:metode-pengujian}
% [ISI DI SINI — Black-Box + SUS]

% ------------------------------------------------------------
\section{DESAIN SISTEM}
\label{sec:desain-sistem}
% ------------------------------------------------------------
% TODO Phase 2: Sub-sections detail scope Can
% Estimasi: ~12-15 halaman total

\subsection{Kebutuhan Pengguna dan Use Case}
\label{subsec:kebutuhan-pengguna}
% TODO Phase 2: Tabel kebutuhan pengguna (POLosan, format Izzah)
% TODO: Use case diagram dengan 3 aktor: Siswa, Admin/Guru, SUPERADMIN (pivot #1)
% Estimasi: ~2 halaman
% [ISI DI SINI + Tabel 3.1 Kebutuhan Pengguna]

\subsection{Alur Pengguna (User Flow)}
\label{subsec:alur-pengguna}
% TODO Phase 2: User flow placement — CAN ISI SENDIRI
% Estimasi: ~1.5 halaman
% [ISI DI SINI + Gambar 3.x User Flow]

\subsection{Desain Level}
\label{subsec:desain-level}
% TODO Phase 2: 3 level — Trust Building, Ambiguous Choice, Risk & Override
% Estimasi: ~3 halaman
% Format: Tabel spesifikasi teknis per level (POLosan)

\subsubsection{Level 1 --- Guided Recognition (Trust Building)}
% [ISI DI SINI + Tabel 3.x Spesifikasi Level 1]

\subsubsection{Level 2 --- Ambiguous Choice (Doubt \& Calibration)}
% [ISI DI SINI + Tabel 3.x Spesifikasi Level 2]

\subsubsection{Level 3 --- Risk \& Override (Human Agency)}
% [ISI DI SINI + Tabel 3.x Spesifikasi Level 3]

\subsection{Mekanisme Human-in-the-Loop (Top-6 Check)}
\label{subsec:hitl-mekanisme}
% TODO Phase 2: REVISE dari existing — Top-6 check (PIVOT #2 16/6/26)
% BUKAN Override Budget lagi
% Estimasi: ~2 halaman
% Konten: Top-3 tampil, Top-6 cek backend, override flow, anti-cheat
% [ISI DI SINI + Gambar 3.x Top-6 Check Flow]

\subsection{Maskot Momo}
\label{subsec:maskot-momo}
% TODO Phase 2: ADD section baru — deskripsi tekstual konsep maskot
% Estimasi: ~1 halaman
% Pilih dari OPSI E "Doodle Cube" atau OPSI G "Eraser Ghost"
% Konten WAJIB: konsep, 5 state ekspresi, prinsip 8-bit geometric,
% text bubble only, NO voice/NLP, NO legs/wings, float only
% [ISI DI SINI + Gambar 3.x Placeholder Sketsa Momo]

\subsection{Wireframe Antarmuka}
\label{subsec:wireframe}
% TODO Phase 2: 4-6 wireframe dengan keterangan fungsi
% Estimasi: ~2.5 halaman
% Minimum: Onboarding, Main Gameplay, Probe UI, Result & Feedback

\subsubsection{Wireframe Onboarding}
% [ISI DI SINI + Gambar 3.x Wireframe Onboarding]

\subsubsection{Wireframe Main Gameplay}
% [ISI DI SINI + Gambar 3.x Wireframe Main Gameplay]

\subsubsection{Wireframe Probe UI (Top-3 + Top-6 Check)}
% [ISI DI SINI + Gambar 3.x Wireframe Probe UI]

\subsubsection{Wireframe Result dan Feedback}
% [ISI DI SINI + Gambar 3.x Wireframe Result]

\subsection{Rancangan Pengujian}
\label{subsec:rancangan-pengujian}
% TODO Phase 2: Black-Box + SUS questionnaire + prosedur pengujian
% Estimasi: ~1.5 halaman
% [ISI DI SINI]

\subsection{Pembagian Tugas dan Alur Pengerjaan}
\label{subsec:pembagian-tugas}
% TODO Phase 2: DUO PATTERN (dari Daffa 3.2.2)
% Estimasi: ~1 halaman
% Diagram pembagian tugas (BIRU=Can, ABU=Dias)
% Narasi: penulis (Can) ↔ rekan (Dias)
% [ISI DI SINI + Gambar 3.x Pembagian Tugas]

% ------------------------------------------------------------
\section{JADWAL PENELITIAN}
\label{sec:jadwal-penelitian}
% ------------------------------------------------------------
% TODO Phase 2: Tabel timeline SINKRON dengan Dias
% Estimasi: ~1 halaman
% Format: 1 tabel shared, kolom Bulan 1-12, rows kegiatan + inisial C/D
% Lihat Style_Guide_v1.md Section 7.3 untuk daftar kegiatan
%
% [ISI DI SINI + Tabel 3.x Timeline]

% ============================================================
% DAFTAR PUSTAKA
% ============================================================
\chapter*{DAFTAR PUSTAKA}
\addcontentsline{toc}{chapter}{DAFTAR PUSTAKA}

% TODO Phase 2: Daftar Pustaka format PENS-IEEE
% Urut by nomor [1], [2], ... sesuai urutan muncul
% Format: Author. (Year). Title. Venue, vol. X, no. Y, pp. Z-W. doi: ...
%
% Referensi Can (22 total — lihat proposal-can.tex existing):
% [1] Ng et al. - AI literacy
% [2] Casal-Otero et al. - AI literacy K-12
% [3] Khosravi et al. - XAI in education
% [4] Mosqueira-Rey et al. - HITL ML state of the art
% [5] Memarian & Doleck - HITL in AIED
% [6] Videnovik et al. - GBL CS education
% [7] Chan, Wan & King - Flow Theory
% [8] Schroeder et al. - Pedagogical agents
% [9] Karran et al. - Confidence visualization
% [10] Pressman - Software Engineering
% [11] Sommerville - Software Engineering
% [12] Sugiyono - R&D
% [13] Creswell - Mixed Methods
% [14] Ishikawa - Quality Control
% [15] Brooke - SUS
% [16] Bangor et al. - SUS empirical
% ... + tambahan Waterfall SDLC reference
% [ISI DI SINI]

\end{document}
